Fix the ordinal $\alpha := \omega^{\omega^2}$.
Let $[\alpha]^2 := \{\{x,y\}: x,y<\alpha,\ x\neq y\}$.
A \,\emph{red/blue edge-colouring} is a map $c:[\alpha]^2\to\{\text{red},\text{blue}\}$.
Write $K_\alpha$ for the complete graph on vertex set $\alpha$ with edge-set $[\alpha]^2$.

A \,\emph{red $K_\alpha$} means: there exists a subset $X\subseteq \alpha$ such that
(1) $(X,<)$ is order-isomorphic to $(\alpha,<)$ (so $X$ has order-type $\alpha$), and
(2) for all distinct $x,y\in X$, $c(\{x,y\})=\text{red}$.

A \,\emph{blue $K_3$} means: there exist $x<y<z<\alpha$ with
$c(\{x,y\})=c(\{x,z\})=c(\{y,z\})=\text{blue}$.

Question: Is it true that every $c:[\alpha]^2\to\{\text{red},\text{blue}\}$ produces either a red $K_\alpha$ or a blue $K_3$?
(Equivalently: does $\alpha\to(\alpha,3)^2$ hold for $\alpha=\omega^{\omega^2}$?)
